\section{A uniform decoder}\label{sec:decoding}

Theorem~\ref{thm:separation} supplies a test for the first incorrect bit of
a prefix: no moment system fixing that prefix can reproduce all the
low-order trace statistics.  To turn this statement into an algorithm,
we must also recognize a correct prefix in bounded time.  The useful
additional fact is that a correct prefix admits a moment system with
strict slack, both in its statistical constraints and in its positivity
constraint.  We first exhibit that slack, then give a rational procedure
which is guaranteed to find a feasible system whenever such slack exists.

\subsection{Statistical constraints for a prefix}

Assume for now that \(p<1\).  Choose an integer \(D\) and a positive dyadic
number \(\zeta=2^{-b_\zeta}\) such that
\begin{equation}\label{eq:decoder-parameters}
 D_*\le D\le C Y^2,\qquad
 \zeta\le \tfrac14,\qquad
 4\zeta<e^{-C_{\mathrm{sep}}B(n,p)},\qquad b_\zeta=O(B(n,p)),
\end{equation}
where \(C_{\mathrm{sep}}\) is the constant in
Theorem~\ref{thm:separation}.  Integer formulas for these choices will be
given below.  Let
\[
 \mathcal I_D=\{I\subseteq\{1,\ldots,n\}:1\le |I|\le D\},
 \qquad N=|\mathcal I_D|,
 \qquad T=512(1+N)\zeta^{-2}.
\]
The number \(T\) is an integer.  From \(T\) independent traces, form the
empirical mean \(\widehat t_I\) of the rank observable defining \(T_I\),
for every \(I\in\mathcal I_D\).  These are observables of the ordinary
trace: if a requested rank is absent, its observable is zero.

For each fixed input \(x\), let \(\mathcal G\) be the event
\begin{equation}\label{eq:empirical-good}
 |\widehat t_I-T_I(x)|\le \zeta/8
 \quad\text{for every }I\in\mathcal I_D.
\end{equation}
Each observable takes values in \(\{0,1\}\), so its empirical mean has
variance at most \(1/(4T)\).  Chebyshev's inequality and a union bound give
\begin{equation}\label{eq:good-probability}
 \mathbb P_x(\mathcal G^c)
 \le \frac{16N}{T\zeta^2}\le\frac1{32}.
\end{equation}
The same event will control all the prefix tests, even though the
prefixes are chosen using the data.

For a proposed prefix, use the remaining sign variables and their moment
matrix \(M(m)=\mathcal M_D(m)\) from \eqref{eq:moment-matrix}.  Its nonempty moment
coordinates form a real vector \(m\).  Consider the rational constraints
\begin{equation}\label{eq:prefix-feasibility}
 m\in[-1,1]^l,\qquad M(m)\succeq0,\qquad
 |\EE_m[T_I]-\widehat t_I|\le\zeta
 \quad(I\in\mathcal I_D),
\end{equation}
where \(l\) is the number of moment coordinates.  Substituting the prefix
and writing each remaining bit as one half of one plus its sign makes
\(\EE_m[T_I]\) an affine function of \(m\).  Its nonconstant coefficients
have absolute sum at most \(1\), by the coefficient bound for \(T_I\)
in Section~\ref{sec:setup}.

\begin{lemma}\label{lem:prefix-slack}
On \(\mathcal G\), a prefix whose first incorrect bit is at its last
position has no solution to \eqref{eq:prefix-feasibility}.  Every correct
prefix has a moment vector \(m^*\in[-1,1]^l\) such that, writing \(G\)
for the dimension of its moment matrix,
\begin{equation}\label{eq:prefix-slack}
 M(m^*)\succeq(\zeta/8)\operatorname{Id}_G,\qquad
 |\EE_{m^*}[T_I]-\widehat t_I|\le\zeta/2
 \quad(I\in\mathcal I_D).
\end{equation}
\end{lemma}

\begin{proof}
If a wrong prefix had a feasible moment vector, then
\eqref{eq:empirical-good} and \eqref{eq:prefix-feasibility} would give
\[
 |\EE_m[T_I]-T_I(x)|\le 9\zeta/8
 \quad(I\in\mathcal I_D).
\]
This contradicts Theorem~\ref{thm:separation} and
\eqref{eq:decoder-parameters}.

For a correct prefix, take the moments of the true completion with
weight \(1-\zeta/8\), and the moments of independent uniform signs on
the remaining positions with weight \(\zeta/8\).  The moment matrix of
the latter distribution is the identity: its \((I,J)\) entry is zero
unless \(I=J\).  The mixture therefore has the stated matrix lower
bound, and its coordinates lie in \([-1,1]\).  Both constituent
distributions give channel expectations in \([0,1]\), so mixing changes
each \(T_I(x)\) by at most \(\zeta/8\).  Equation~\eqref{eq:empirical-good}
then bounds the discrepancy from \(\widehat t_I\) by \(\zeta/4\),
which is stronger than required.
\end{proof}

\subsection{Finding a moment system with slack}

The next lemma gives the bounded procedure needed for
\eqref{eq:prefix-feasibility}.  It only promises to succeed when there
is slack; every reported solution satisfies the exact inequalities.
This distinction avoids any need to decide arbitrary instances of
semidefinite feasibility.
The procedure is a projected relaxation method for inequalities, in
the tradition of Agmon and Motzkin--Schoenberg
\cite{Agmon1954,MotzkinSchoenberg1954}.  We give the iteration and its
rounding analysis explicitly.

\begin{lemma}[Rational feasibility with slack]\label{lem:rational-feasibility}
Let \(r\ge0\) and \(D\ge1\) be integers, and let \(m\) consist of the
\[
 l=\sum_{j=1}^{\min(2D,r)}\binom rj
\]
nonempty sign moments.  Let \(M(m)\) be their moment matrix, of
dimension
\[
 G=\sum_{j=0}^{\min(D,r)}\binom rj.
\]
Suppose a finite list of rational affine inequalities
\(\ell_t(m)\le b_t\) is given, each with nonconstant coefficient vector
of \(\ell^1\)-norm at most \(1\), together with a dyadic
\(0<\zeta\le1/4\).

There is a deterministic rational procedure which either returns a
vector satisfying
\[
 m\in[-1,1]^l,\qquad M(m)\succeq0,\qquad \ell_t(m)\le b_t
 \quad\text{for all }t,
\]
or reports failure.  It returns a vector whenever some
\(m^*\in[-1,1]^l\) satisfies
\[
 M(m^*)\succeq(\zeta/8)\operatorname{Id}_G,\qquad
 \ell_t(m^*)\le b_t-\zeta/2\quad\text{for all }t.
\]
Put \(\lambda=\zeta/(16G)\).  The procedure performs at most
\(\lceil10(1+l)/\lambda^2\rceil\) feasibility checks, each taking a
number of bit operations polynomial in the list dimensions, the
input bit lengths, and \(\log(1/\zeta)\).  These time bounds hold
without the slack assumption.
\end{lemma}

\begin{proof}
If \(l=0\), the matrix is \([1]\) and all inequalities are constant,
so direct rational comparisons suffice.  Suppose \(l>0\).  We describe
exact separating directions, then an iteration with bounded precision.

\paragraph{Directions from violated constraints.}
At a current point \(m^{\mathrm{cur}}\in[-1,1]^l\), first check the
affine inequalities and the matrix condition exactly.  If all hold,
return the current vector.  A violated inequality
\(\ell_t(m^{\mathrm{cur}})>b_t\) gives the direction \(a\) equal to
minus its nonconstant coefficient vector.  Then \(\norm{a}_2\le1\);
under the slack assumption,
\[
 \langle m^*-m^{\mathrm{cur}},a\rangle\ge\zeta/2\ge\lambda.
\]

If instead the matrix condition fails, find a nonzero rational vector
\(v\) with \(v^\top M(m^{\mathrm{cur}})v<0\).  The exact construction
of \(v\) is given just below.  Write the associated affine form as
\[
 v^\top M(m)v=c+h\cdot m,\qquad
 a=\frac{h}{Gv^\top v}.
\]
Each matrix entry is either the constant \(1\) or one moment
coordinate.  Consequently
\[
 \norm{h}_1\le\norm{v}_1^2\le G\norm{v}_2^2,
 \qquad \norm{a}_2\le1.
\]
The slack matrix lower bound gives
\[
 \langle m^*-m^{\mathrm{cur}},a\rangle
 =\frac{v^\top(M(m^*)-M(m^{\mathrm{cur}}))v}{Gv^\top v}
 \ge\frac{\zeta}{8G}\ge\lambda.
\]
Thus either violated condition provides a rational direction satisfying
\begin{equation}\label{eq:feasibility-direction}
 \norm{a}_2\le1,\qquad
 \langle m^*-m^{\mathrm{cur}},a\rangle\ge\lambda
\end{equation}
whenever a slack solution exists.  If a constant affine inequality is
violated, its direction is zero; this cannot happen in the slack case,
and in all other cases the iteration is still bounded.

\paragraph{An exact matrix check.}
For a rational symmetric matrix, eliminate a positive diagonal pivot
by taking its Schur complement.  Congruence by an invertible triangular
matrix shows that the original matrix is positive semidefinite exactly
when the positive pivot and this Schur complement are.  A negative
diagonal supplies a negative unit vector.  If a remaining diagonal at
\(i\) is zero and its entry \(b\) in column \(j\) is nonzero, let \(c\)
be the diagonal at \(j\).  We may assume \(c\ge0\), since otherwise
there is already a negative diagonal.  The rational vector
\[
 y=e_i-\frac{b}{1+c+|b|}e_j
\]
has quadratic form
\[
 \frac{b^2}{(1+c+|b|)^2}
 \bigl[-2(1+c+|b|)+c\bigr]<0.
\]
If no negative vector is found, continue with a positive diagonal
pivot; when none remain, all entries are zero and the remainder is
positive semidefinite.

For completeness, a negative vector of a Schur complement lifts to
one of the original matrix using only rational arithmetic.  If \(P\)
is the original block of eliminated pivots and \(B\) its block of
entries against the remaining indices, the lift of \(y\) is
\[
 v=(-P^{-1}By,y).
\]
Its quadratic form is exactly the quadratic form of \(y\) for the
Schur complement.

These operations have polynomial bit complexity.  Indeed, clear
denominators in a rational matrix with \(t\)-bit entries; even the
product of all denominators has at most \(G^2t\) bits.  If \(H\) is
the resulting integer matrix and its entries have at most \(s\) bits,
the determinant expansion bounds every \(j\)-row minor by
\(j!\,2^{js}\).  Its bit length is therefore \(O(G(s+\log G))\).
Schur-complement entries are ratios of minors of the original matrix,
and the inverse formula for \(P\) gives the same bound for the lift.
Fraction-reduced elimination, with rational comparisons at each step,
therefore computes the check and the negative vector in polynomially
many bit operations.  This argument applies however small a positive
pivot or a negative quadratic form may be.

\paragraph{Iteration and rounding.}
Start with \(m^{\mathrm{cur}}=0\).  Choose the dyadic grid with mesh
\[
 \eta=2^{-g},\qquad
 g=\left\lceil\log_2\frac{100(1+l)}{\lambda^2}\right\rceil.
\]
The integer \(g\) is computed by rational comparisons with powers of
two; thus \(\eta\le\lambda^2/(100(1+l))\), with
\(O(1+\log((1+l)/\lambda))\) fractional bits.  After an
unsuccessful check, add \(\lambda a\), clip every coordinate to
\([-1,1]\), and round to the nearest point of this grid in the box.
The endpoints of the box belong to the grid.  Repeat up to the
stated number of checks, reporting failure if none accepts.

Under the slack assumption, \eqref{eq:feasibility-direction} implies
\[
 \norm{m^{\mathrm{cur}}+\lambda a-m^*}_2^2
 \le \norm{m^{\mathrm{cur}}-m^*}_2^2-\lambda^2.
\]
Clipping cannot increase the distance to \(m^*\).  Since both the
clipped vector and \(m^*\) lie in the box, rounding changes the squared
distance by at most \(4l\eta+l\eta^2<\lambda^2/2\).  Each failed check
thus decreases squared distance by at least \(\lambda^2/2\).
Its initial value is at most \(l\), so acceptance must occur before
the cap.

Finally, rounding keeps every iterate on a single grid of the indicated
bit length.  The exact matrix check, the affine comparisons, the
formation and normalization of \(h\), and one update involve
polynomially many rational operations on numbers of polynomial bit
length.  Temporary denominators within an update therefore have
polynomial bit length, and they do not accumulate across updates.
The fixed cap and these bounds apply to every rational input list,
whether or not it admits a slack solution.
\end{proof}

For a prefix test, write each absolute-value constraint in
\eqref{eq:prefix-feasibility} as two affine inequalities and apply
Lemma~\ref{lem:rational-feasibility}.  Lemma~\ref{lem:prefix-slack}
verifies exactly the slack hypotheses for a correct prefix.  Every
acceptance, on any data, gives an actual moment system satisfying
\eqref{eq:prefix-feasibility}.

The decoder starts with the empty prefix and tests both one-bit
extensions.  If precisely one accepts, it appends that bit; otherwise
it appends \(0\).  It repeats until the prefix has length \(n\).
The latter rule specifies a length-\(n\) output even for data on which
the statistical estimates fail.  There are at most \(2n\) tests.

\subsection{Integer parameters and bit complexity}

We now make the choices in \eqref{eq:decoder-parameters} using only
the binary input.  This also accounts for the size of every list and
every rational calculation in the decoder.  Continue to assume
\(p<1\), and use \(q_*,Q,Y\) from \eqref{eq:aux-parameters}.  Define
the positive integers
\begin{align*}
 u&=1+\lceil\log_2 n\rceil,
 &v&=1+\lfloor\log_2(1/q_*)\rfloor,\\
 y&=1+\lceil u/v\rceil,
 &w&=1+\lceil\log_2 y\rceil,
 &a_p&=\lceil1/p\rceil.
\end{align*}
All logarithms appearing in these integer formulas mean integer
length comparisons.  In particular, the integer part of a logarithm
of a rational number is found by comparing its numerator with
power-of-two multiples of its denominator.  The rational maximum
defining \(q_*\) is also found by an integer comparison.

The elementary comparisons
\[
 u\asymp\log n,\qquad v\asymp Q,\qquad
 Y\le y\le C Y,\qquad
 w\asymp1+\log Y,\qquad p^{-1}\le a_p\le2p^{-1}
\]
have absolute constants.  For the asserted lower bound on \(y\),
use \(u\ge(\log n)/\log2\) and \(v\le Q/\log2\).
Together with \(Y\asymp\mathsf X(n,p)\), these comparisons give
\begin{equation}\label{eq:integer-budget-proxy}
 B(n,p)\le a_pu y^2w^6\le C B(n,p).
\end{equation}
Choose fixed positive integers \(K_D,K_m\), sufficiently large for
the constants in Theorem~\ref{thm:separation}, and set
\begin{equation}\label{eq:integer-decoder-parameters}
 D=K_Dy^2,\qquad b_\zeta=K_m a_pu y^2w^6,\qquad
 \zeta=2^{-b_\zeta}.
\end{equation}
Since \(D_*\le C Y^2\) and \(B(n,p)\ge\log2\), these choices ensure
all of \eqref{eq:decoder-parameters}.  The integers \(D,b_\zeta\), in binary,
are computable in time polynomial in the input bit lengths.  Writing
the denominator \(2^{b_\zeta}\) subsequently costs \(O(b_\zeta)\) bits, which will
be included in the bounds below.

We next describe how to form the rational constraints.  For
\(I=\{l_1<\cdots<l_k\}\) and \(i_1<\cdots<i_k\), the coefficient
of \(X_{i_1}\cdots X_{i_k}\) in \(T_I\) is
\begin{equation}\label{eq:explicit-channel-coefficient}
 \binom{i_1-1}{l_1-1}
 \prod_{j=1}^{k-1}
 \binom{i_{j+1}-i_j-1}{l_{j+1}-l_j-1}
 \,p^{l_k}q^{i_k-l_k},
\end{equation}
when all the gap counts are possible, and is zero otherwise.
Indeed, the selected positions must be retained, and the prefix and
successive open gaps must contain the indicated numbers of additional
retentions.  Positions after \(i_k\) are unrestricted.  Formula
\eqref{eq:explicit-channel-coefficient} follows by multiplying these
independent binomial probabilities.

If the input representation is \(p=a/b\), every denominator in
\eqref{eq:explicit-channel-coefficient} divides \(b^n\).
Binomial coefficients have \(O(n)\) bits.  Sign expansion adds a
denominator dividing \(2^D\), and empirical means have denominator
\(T\).  Consequently the constraint coefficients have polynomial
bit length in \(n+L+D+\log T\).  Enumerating the rank and input tuples,
then expanding the resulting bit monomials and substituting the
prefix, constructs them directly.  No list of complete trace masks
is required.

All these enumerations fit the permitted resource scale.  From
\(D=O(Y^2)\), \(Y\asymp\mathsf X(n,p)\), and \eqref{eq:budget},
\begin{equation}\label{eq:decoder-size-bound}
 D\log n=O(B(n,p)).
\end{equation}
The rank list, every matrix basis, and every moment list have size
at most \((n+1)^{2D}=\exp(O(B(n,p)))\).  Input tuple lists satisfy the
same bound.  A sign expansion has at most \(2^D\) terms; products
of these enumeration bounds remain \(\exp(O(B(n,p)))\).  Subsets
larger than the available number of positions are simply absent,
so this statement includes \(D>n\) and prefixes with very few
remaining positions.

By \eqref{eq:decoder-parameters}, the trace count \(T\) is
\(\exp(O(B(n,p)))\).  For every prefix, the iteration cap in
Lemma~\ref{lem:rational-feasibility} is bounded by
\[
 2560(1+l)G^2\zeta^{-2}+1=\exp(O(B(n,p))).
\]
The bit length of its grid is
\(O(b_\zeta+\log G+\log(1+l))=O(B(n,p))\).
The preceding coefficient bounds and that lemma therefore bound
all computation by a fixed polynomial in
\begin{equation}\label{eq:decoder-total-resource}
 n+L+\exp(O(B(n,p))).
\end{equation}
In particular, this includes the construction of parameters and
constraints, exact rational arithmetic, and the at most \(2n\)
prefix tests.

\begin{proof}[Proof of Theorem~\ref{thm:main}]
If \(p=1\), request one trace and output it.  For the universal
halting specification, output \(0^n\) if the supplied well-formed
trace has a different length.  The actual channel always supplies
the input itself, and this branch takes polynomially many bit
operations in \(n+L\).

For \(p<1\), use \eqref{eq:integer-decoder-parameters}, request the
specified \(T\) traces, and compute their empirical statistics.
Run the prefix rule specified above.  This defines the machine on
every stream of well-formed traces.

On \(\mathcal G\), induction on the prefix length shows that
precisely the correct extension accepts: the incorrect extension
is infeasible by Lemma~\ref{lem:prefix-slack}, and the correct one
has the slack that guarantees acceptance.  Thus the output is \(x\)
on \(\mathcal G\), whose probability is at least \(31/32>2/3\) for
each fixed \(x\).

Each supplied trace has at most \(n\) bits and one delimiter.
Reading them, and computing each of the finitely many rank
observables, costs polynomially many bit operations in \(n+T+N\);
the empirical counters have \(O(\log T)\) bits.  All loops, including
all feasibility tests, have the input-dependent caps already given.
Thus \eqref{eq:decoder-total-resource} is a bound for every
well-formed trace stream, independently of the success event.
The machine can ignore its internal random bits.

To be explicit about the computational model, the described finite
lists may be stored sequentially on tapes; replacing an indexed
lookup by a scan incurs at most a polynomial overhead in their
total encoded length.  Integer addition, multiplication, division,
and fraction reduction have elementary polynomial-time multitape
implementations.  The bit-length and loop bounds above therefore
give absolute constants \(K,d\) for which the trace count is at most
\(\exp(KB(n,p))\) and the total bit cost is at most
\((n+L+\exp(KB(n,p)))^d\), after enlarging \(K,d\) to absorb fixed
factors.  This enlargement is uniform because \(B(n,p)\ge\log2\).
Choose absolute \(C,c\) large enough, with \(C\ge K\) and \(c\ge d\).
Both the sample bound \(M_C(n,p)\) and the required runtime
\((n+L+M_C(n,p))^c\) follow.  All constants used in parameter
selection are fixed integers in a single finite program.
\end{proof}
